\subsection{Platform Configuration, Policies and Audit}
\label{sec:impl_configuration}

\paragraph{Runtime settings}
Operational limits --- concurrent programme enrolments, career paths per programme, retention windows --- are runtime settings rather than deployment constants, resolved per organisation by \path{abridgeai/core/runtime_settings.py}. The effective value is the first of four levels that supplies one: a \texttt{system\_settings} row for this organisation, a row with a null organisation (the deployment default), the environment variable named by the setting's spec, and finally the code default declared in the settings registry.

The last two levels are what make adding a setting safe. A deployment that has never heard of a new key still resolves it, so a release does not require a configuration change to land first. The resolution API returns every candidate value alongside the effective one and the level it came from, because an administrator asking ``why is this 3?'' needs to see which layer won rather than be told the answer.

One row set per organisation is cached in process under a short TTL, since settings are read on many requests and written rarely.

\paragraph{Policy documents}
\path{abridgeai/features/policies/} holds terms, privacy notices, and academic regulations. Versioning is copy-on-write: a published version is immutable and correcting it opens a new draft, so the exact text a reader agreed to remains retrievable. Audience is a set of roles, and the empty set means public --- including signed-out visitors.

\paragraph{Audit}
Two audit trails run side by side. The typed authentication and access-control trail records sign-in, session, and role-assignment events; the HTTP audit log records mutating requests with actor, method, path, affected row, and a column-level diff, written asynchronously through an actor pool so the request path carries no additional latency.

Both are append-only, and the guarantee is enforced in the database rather than in application code: a trigger rejects \texttt{UPDATE} and \texttt{DELETE} on the audit tables outright, with deletion permitted only inside a narrowly scoped maintenance session used by the retention job. This is the property that gives the trail its value --- a participant cannot remove evidence of what they did, and neither can an administrator.

\diagramplaceholder{fig:impl_settings_resolution}{Runtime setting resolution order}{A flow chart resolving one setting key for one organisation through the four levels in order --- organisation row, deployment-default row, environment variable, code default --- returning at the first hit. Show the in-process cache and its TTL, and show that all candidates are reported alongside the winning value and its source.}

\diagramplaceholder{fig:impl_audit_path}{Audit write path and append-only enforcement}{Show a mutating request producing an audit record dispatched asynchronously to the actor pool, so the response does not wait. Alongside it, show the database trigger rejecting \texttt{UPDATE} and \texttt{DELETE}, and the single scoped exception by which the retention job may delete after the configured window.}
